\section{Introduction}
Eye movements provide a time-resolved behavioral trace of reading. Fixation durations,
refixations, skipping, and go-past measures respond to lexical, contextual, and reader
factors, making natural-reading gaze data a useful setting for studying how linguistic
pressure differs across reader groups \cite{rayner1998eye,duchowski2017eye}.
Danish natural reading is especially relevant because orthography, morphology, and
vocalic boundary patterns create pressures that are not identical to English-centered
benchmarks. Prior Danish dyslexia-reader prediction from eye movements has shown that
reader group can be predicted from gaze behavior, but the interpretability of such
predictions and the linguistic pressures driving them remain underdeveloped.

Language-model predictability offers a bridge between psycholinguistic reading-time
theory and modern NLP. The central question in this paper is not only whether
dyslexia-labeled and typical/control readers can be distinguished, but whether the
distinction is carried by text exposure, global reading speed, or participant-level
sensitivity to DFM predictability. We answer this with a frozen confirmatory analysis:
DFM predictability sensitivity and cross-fitted residualized gaze-cost profiles provide
the main predictive signal.
